\section{Appendix: Identification and upper-bound proofs}\label{sec:appendix-upper}

This appendix collects the analytic ingredients for \cref{obj:prop:identification,obj:thm:observable-upper}. The causal argument connects the observed-law class to admissible full-data completions. The upper-bound analysis uses armwise sensitivity to organize localization, polynomial approximation, factorial estimation, and clipping. Throughout, the estimator and its sampling conventions are those of \cref{obj:def:armwise-estimator}; the detailed arguments appear in \cref{sec:deferred-proofs}.

\subsection{Causal completion and value identification}

The two causal clauses of \cref{obj:prop:identification} have distinct mathematical roles. Existence of a completion ensures that every law in \cref{obj:def:observed-class} has a causal interpretation satisfying \cref{obj:ass:consistency,obj:ass:conditional-exchangeability}. Value identification then establishes that all completions in \cref{obj:def:causal-class} yield the observed value of their marginal law. Together, these clauses connect the observed statistical problem to the causal oracle value under the stated overlap restriction.

Positive overlap supplies both conditional outcome means in every occupied cell. Consistency and conditional exchangeability relate these observed means to the corresponding conditional potential-outcome means. The oracle can choose an arm separately in each cell, so its value takes the cell-mass-weighted maximum appearing in \cref{obj:def:observed-value}. Cells of zero probability contribute zero under the conventions of \cref{obj:prop:identification}.

The final clause of \cref{obj:prop:identification} supplies the cellwise representation used in estimation. At an admissible vector of observed atom probabilities, the overlap-truncated denominators in \cref{obj:def:armwise-extension} agree with the arm masses. The extension therefore represents the target contribution of that cell, including its zero-mass case. Its domain also accommodates the nonnegative vectors generated by pilot localization.

\subsection{Armwise sensitivity}

The local analysis retains the sensitivity of each treatment arm separately. The following bounds control perturbations of the extension and relate the resulting armwise weights to the observed atom masses.

\begin{lemmav}{lem:armwise-modulus}[Armwise perturbation and mass bounds]
Let \(\epsilon\) be the overlap parameter, with \(0<\epsilon\leq 1/2\). For every pair of four-vectors \(u,v\) satisfying
\begin{itemize}
\item \textbf{(Nonnegativity.)} \(u,v\in\mathbb R_+^4\).
\item \textbf{(Nonzero comparison.)} \(v\neq 0\).
\end{itemize}
define the armwise perturbations by
\[
\Delta_a=\sum_{y=0}^1|u_{ay}-v_{ay}|,\qquad a\in\{0,1\}.
\]
With \(F_\epsilon\), \(s\), and \(D_a\) as in \cref{obj:def:armwise-extension},
\[
|F_\epsilon(u)-F_\epsilon(v)|
\leq
2\sum_{a=0}^1\Delta_a
+
2\sum_{a=0}^1\frac{s(v)}{D_a(v)}\Delta_a.
\]

For every observed law \(\mathbb P\) on
\(\{0,\ldots,d-1\}\times\{0,1\}\times\{0,1\}\) and every cell \(x\in\{0,\ldots,d-1\}\), let \(v=q_x\) be its vector of observed atom probabilities and \(p_x\) its cell probability:
\[
v_{ay}=\mathbb P\bigl(\{(x,a,y)\}\bigr),
\qquad
p_x=\sum_{a=0}^1\sum_{y=0}^1v_{ay}.
\]
If
\begin{itemize}
\item \textbf{(Dimension.)} \(d\) is an integer with \(d\geq2\).
\item \textbf{(Observed class.)} \(\mathbb P\in\mathcal D_{d,\epsilon}^{\mathrm{obs}}\), as in \cref{obj:def:observed-class}.
\item \textbf{(Positive cell mass.)} \(p_x>0\).
\end{itemize}
then, with \(s_a\) as in \cref{obj:def:armwise-extension},
\[
\sum_{a=0}^1\frac{p_x}{s_a(v)}
\sum_{y=0}^1\sqrt{v_{ay}}
\leq
2\sqrt{2}\sqrt{\frac{p_x}{\epsilon}}.
\]
\end{lemmav}

The armwise perturbation \leanref{sym:\Delta_a}{\ensuremath{\Delta_a}} measures the total coordinate change within one treatment arm. Nonnegativity preserves the mass interpretation of the extension, while a nonzero comparison vector makes each overlap-truncated denominator positive. The first inequality applies even when the perturbed vector is zero, which is useful for localized estimation near empty cells.

The second inequality concerns occupied cells in the observed-law class. It combines the armwise sensitivity weights with square roots of atom masses, the quantities governing count fluctuations. Its dependence on overlap is proportional to the inverse square root of the overlap floor. This mass-sensitive control is the analytic ingredient supporting the single inverse power of overlap in the squared-risk guarantee of \cref{obj:thm:observable-upper}.

\subsection{Factorial statistics and sampling conventions}

Polynomial estimation requires statistics corresponding to powers of atom probabilities. A deterministic split gives a multinomial construction, which we record before discussing its relationship to the active estimator's Poisson construction.

\begin{definitionv}{def:fixed-sample-factorial}[Centered factorial statistic \(U_{\boldsymbol h}\)]
The centered multinomial factorial statistic is
\[
U_{\boldsymbol h}(N^{\mathrm e};b)
=
\sum_{\boldsymbol t\le\boldsymbol h}
\left\{
\prod_{j\in\mathcal J_d}
\binom{h_j}{t_j}(-b_j)^{h_j-t_j}
\right\}
\frac{\prod_{j\in\mathcal J_d}(N_j^{\mathrm e})_{t_j}}
{(n_{\mathrm e})_{|\boldsymbol t|}},
\qquad
|\boldsymbol h|\le n_{\mathrm e}.
\]
Here \(\mathcal J_d\) is the observed atom index set, \(b=(b_j)_{j\in\mathcal J_d}\) is the centering vector, and \(\boldsymbol h,\boldsymbol t\) are nonnegative integer multi-indices, with componentwise ordering and total orders
\[
|\boldsymbol h|=\sum_{j\in\mathcal J_d}h_j,
\qquad
|\boldsymbol t|=\sum_{j\in\mathcal J_d}t_j.
\]
The deterministic sample split and its atom counts are
\[
n_{\mathrm p}=\lfloor n/2\rfloor,
\qquad
n_{\mathrm e}=n-n_{\mathrm p},
\qquad
I_{\mathrm p}=\{1,\ldots,n_{\mathrm p}\},
\qquad
I_{\mathrm e}=\{n_{\mathrm p}+1,\ldots,n\},
\]
\[
N_j^{\mathrm p}
=\sum_{i\in I_{\mathrm p}}\mathbf 1\{O_i=j\},
\qquad
N_j^{\mathrm e}
=\sum_{i\in I_{\mathrm e}}\mathbf 1\{O_i=j\},
\qquad
N^{\mathrm e}=(N_j^{\mathrm e})_{j\in\mathcal J_d}.
\]
The falling factorials are
\[
(m)_k=m(m-1)\cdots(m-k+1),
\qquad
(m)_0=1.
\]
\end{definitionv}

The centered multinomial statistic \leanref{sym:U_{\boldsymbol h}(N^{\mathrm e};b)}{\ensuremath{U_{\boldsymbol h}(N^{\mathrm e};b)}} uses the joint factorial moments of a fixed-size evaluation histogram. Its denominator depends on the total order of each term, reflecting the common multinomial sample size. The order restriction ensures that these normalizations are defined. Under \cref{obj:ass:iid-sampling}, the two deterministic subsamples are independent, allowing the centering vector to be selected from the pilot sample.

The active branch of \cref{obj:def:armwise-estimator} uses a separate sampling convention: an auxiliary Poisson count and fair marks generate pilot and evaluation counts. In the uncapped Poisson experiment, the evaluation atom counts are independent Poisson variables. The coordinatewise statistic defined there consequently normalizes each falling factorial by a power of the evaluation intensity. Both constructions translate centered polynomial terms into count statistics, with their respective multinomial and Poisson normalizations. The risk guarantee in \cref{obj:thm:observable-upper} concerns the capped Poisson construction specified in \cref{obj:def:armwise-estimator}.

\subsection{Risk analysis by estimator branch}

The branches of \cref{obj:def:armwise-estimator} separate the roles of empirical estimation, bounded loss, and local polynomial estimation. For a bounded alphabet, the relevant quantities are empirical cell weights and arm-specific outcome means, including the prescribed zero-count convention. Since the alphabet cutoff is universal, the logarithmic alphabet scale remains bounded in this branch. Its contribution to \cref{obj:thm:observable-upper} is therefore assessed within a fixed range of alphabet sizes.

For larger alphabets in the saturated regime, the estimator returns the midpoint of the unit interval. The target lies in that interval by \cref{obj:def:observed-value}, giving a squared-error bound of one quarter. This branch supplies the constant-risk part of \cref{obj:thm:observable-upper}.

The active branch assigns separate mathematical roles to localization, approximation, factorial variance, and clipping. Pilot localization determines a rectangle for each cell's atom probabilities. Its radii include both a mass-dependent fluctuation term and a positive additive term, so the polynomial coordinates remain defined when pilot counts vanish. The perturbation bound in \cref{obj:lem:armwise-modulus} measures variation over these rectangles through the armwise widths and sensitivities.

Polynomial bias concerns the discrepancy between the extension and its tensor Jackson approximation on the pilot rectangle. The approximation degree grows with the logarithmic alphabet scale, while localization determines the region on which the approximation is used. Factorial estimation then represents the polynomial through evaluation counts. Independence in the uncapped Poisson experiment separates the pilot choice of polynomial from evaluation noise; factorial moments govern the variance cost of estimating its centered terms. The tuning constants in \cref{obj:thm:observable-upper} calibrate these approximation and variance roles jointly.

Clipping uses the same armwise geometry. The interval around the center's contribution is scaled by the armwise localization widths and sensitivity weights specified in \cref{obj:def:armwise-estimator}. It bounds each polynomial statistic before aggregation. Projection of the aggregate onto the unit interval aligns its range with that of the target and cannot increase squared error for a target in that interval.

Finally, the Poisson cap connects the auxiliary experiment to the available sample of exactly \(n\) observations. On the cap event, the statistic uses the selected initial observations and their marks; overflow receives the prescribed value zero. The bounded range permits the overflow probability to enter the risk accounting as a controlled discrepancy. Conditional averaging over the auxiliary count and marks then produces the deterministic observed-data estimator. With the proved calibration \(D_0=2\) and the admissible range \(d\ge2\), \cref{obj:thm:observable-upper} supplies the uniform risk bound through the saturated and active branches. The empirical branch belongs to the general parameterized definition and is inactive under this calibration.
